\begin{lemma}[Claim 4.6: the initial Case 2 coloring]
Fix a certified ETT and its support class
\[
 L_1=\noderef{ExtensionBaseSupportClass}(v_n),
\]
with exact final slot \(n\) (so its successor index is the number of
rounds), an exact earlier member
\(n^-\in L_1\setminus\{n\}\), \(n^-<n\), and
\[
 \delta_{n^-}=\gamma_n,\qquad
 \gamma(\text{slot }n)=\gamma_n,\qquad
 \delta(\text{slot }n)=\delta_n.                              \tag{4}
\]
Let \(T_j,T_n,T_j(v_n),T'_j\) denote, respectively, the reference tree,
terminal tree, exact prefix seed, and Case 2 closure.  Require the ordered
initial-segment relations
\[
 T_j(v_n)\preccurlyeq T_j\preccurlyeq T_n,\qquad
 T_j(v_n)\preccurlyeq T'_j,
\]
and require the terminal displayed vertex of the seed \(T_j(v_n)\) to be
exactly \(v_n\).  Let
\(\varphi_{j-1},\varphi_n,\sigma_n,\varrho_1\) be the displayed
colorings; assume the last three are proper and leave exactly \(e_0\)
uncolored.  Let \(D_{j-1}\), \(D\), \(R\), and \(S'\) be the finite
previous control set, terminal control set, restoration-color set, and
support-color set, with the exact identity
\[
 S'=\bigcup_{t\in L_1}\mathcal T.\mathrm{stored}(t).
\]

Retain all of the following Case 2 data explicitly.
\begin{enumerate}
\item \(T'_j\) is a \(\varrho_1\)-\noderef{TashkinovClosure} of
      \(T_j(v_n)\), its vertex set is elementary, and
      \[
       \overline{\varrho}_1(v)\cap S'=\varnothing
       \ (v\in V(T'_j)\setminus\{v_n\}),\qquad
       \overline{\varrho}_1(v_n)\cap S'=\{\delta_n\}.          \tag{C2}
      \]
      The colors \(\gamma_n,\delta_n,\gamma_{1_1}\) belong to \(S'\),
      and \(\gamma_n\ne\delta_n\).
\item \(\sigma_n\) is \((T_n,D,\varphi_n)\)-stable,
      \(\gamma_n,\delta_n\in D\), and Claim 4.3(b3) has the exact form
      \[
       \varrho_1(a)\ne\sigma_n(a)\Longrightarrow
       a\notin E(G[V(T_j)])\text{ and }
       \sigma_n(a)=c,\ \varrho_1(a)=d
       \text{ for some }c,d\in R,                             \tag{b3}
      \]
      where \(\gamma_n,\delta_n\notin R\).
\item The terminal PE boundary facts are
      \[
       \partial_{\varphi_n,\gamma_n}(T_n)=\{f_n\},
      \]
      every \(\delta_n\)-colored boundary identity of \(T_n\) is
      incident with \(T_j(v_n)\), \(f_n\) is incident with \(v_n\), and
      \(\delta_n\in\overline{\varrho}_1(v_n)\).
\item The entire \((\gamma_n,\delta_n)\)-component through \(v_n\)
      under \(\varrho_1\) is represented by a specified simple path
      \(P\) from \(v_n\) to a vertex
      \(x\notin V(T_n)\cup V(T'_j)\).  The path contains a vertex of
      \(T_n\) other than \(v_n\).
\item Claim 4.3(b1) is stated directly:
      \(\varrho_1\) is
      \((T_j(v_n)-v_n,D_{j-1},\varphi_{j-1})\)-stable.
      Its b2 data are the vertexwise equality
      \[
       \overline{\varrho}_1(v)=\overline\varphi_{j-1}(v)
       \quad(v\in V(T_j)\setminus\{v_n\}),                    \tag{b2v}
      \]
      the inclusion
      \[
       \overline\varphi_{j-1}(v_n)
       \subseteq\overline{\varrho}_1(v_n)\cup\{\gamma_{1_1}\},
      \]
      and exact equality of the \(\varrho_1\)- and
      \(\varphi_{j-1}\)-colors on every displayed edge of \(T_j\)
      appended at or after \(T_j(v_n)\).
      The whole \(T_j\) is a
      \(\varphi_{j-1}\)-Tashkinov augmentation sequence from
      \(T_j(v_n)\), and
      \(\delta_n\notin\overline\varphi_{j-1}(T_j)\).
\item Claim 4.3(b4) is also explicit: for every exact
      \(t\in L_1\setminus\{n\}\), whenever its certified fields are
      \(\gamma_t,\delta_t\), these colors are \(T_t\)-strongly
      interchangeable under \(\varrho_1\).  Both colors belong to
      \(S'\), neither equals \(\delta_n\), and the support-order premise
      says that every vertex-node of every component occurring in that
      strong-interchangeability assertion lies in \(V(T_n)\).
\end{enumerate}

Then there are a vertex \(u\in V(T'_j)\setminus V(T_n)\) and an exact
simple exit suffix \(Q=P[x,u]\) with at least one edge identity such that
\(u\) is the only vertex-node of \(Q\) in
\(V(T_n)\cup V(T'_j)\).  For every
\(\beta\in\overline{\varrho}_1(u)\), there are a vertex \(z\in V(T_j)\)
and its exact position \(q_z\ge |T_j(v_n)|\), chosen minimally among
later vertices missing \(\beta\), or terminally when there is no such
vertex, for which the following holds.  Define
\[
 \varrho_0=
 \noderef{OutsideTreeColorSwap}(T'_j,\varrho_1,\beta,\delta_n),
 \qquad
 \mu_0=
 \noderef{TwoColorComponentSwap}
   (G,\varrho_0,\gamma_n,\beta,\mathsf{inl}(u)).               \tag{6}
\]
The coloring \(\mu_0\) is proper, leaves exactly \(e_0\) uncolored, and
satisfies:
\begin{enumerate}
\item[(e1)] \(\mu_0\) is
 \((T_j(v_n)-v_n,D_{j-1},\varphi_{j-1})\)-stable.
\item[(e2)] For every
 \[
  v\in V(T_j(z)\cup T'_j(u))\setminus\{u,z\},
 \]
 one has
 \(\overline{\mu}_0(v)=\overline{\varrho}_1(v)\).  Every exact displayed
 edge identity in \(T_j(z)\cup T'_j(u)\) has the same \(\mu_0\)- and
 \(\varrho_1\)-color.  Moreover,
 \[
  \beta\in\overline{\varrho}_1(z)
  \Longrightarrow \delta_n\in\overline{\mu}_0(z).
 \]
 Here \(T_j(z)\) consists of positions through \(q_z\), and the seed
 edges of this prefix are also present in the exact prefix \(T'_j(u)\).
\item[(e3)] The exact color identity in (4) is retained, and
 \[
  \gamma_n=\delta_{n^-}\in\overline{\mu}_0(u),
  \qquad \beta\notin\overline{\mu}_0(u).
 \]
\item[(e4)] For every exact \(t\in L_1\setminus\{n\}\), its certified
 colors \(\gamma_t,\delta_t\) are \(T_t\)-strongly interchangeable under
 \(\mu_0\).
\item[(e5)]
 \[
 \overline{\mu}_0(T'_j-u)\cap S'
 =\overline{\mu}_0(v_n)\cap S'=\{\delta_n\},
 \]
 and, for every exact edge position \(a\ge |T_j(v_n)|\) of \(T'_j\)
 and every \(c\in S'\setminus\{\delta_n\}\),
 \[
                  \mu_0(T'_j.\operatorname{edge}(a))\ne c .
 \]
\end{enumerate}
\end{lemma}
\begin{proof}
Apply \noderef{ExtensionBaseCaseTwoObservation20} to the explicit
terminal stability, the color-control projection of b3, color-exclusion,
boundary-location, and closure premises.  Here the required seed inclusion
in \(T_n\) is the composite of the two ordered initial-segment relations,
using \noderef{TreeSequenceInitialSegment}; its induced vertex-set
inclusion is obtained from \noderef{TreeSequenceVertexSet}.  It gives
\[
 \partial_{\varrho_1,\gamma_n}(T_n)=\{f_n\},\qquad
 \partial_{\varrho_1,\delta_n}(V(T_n)\cup V(T'_j))=\varnothing. \tag{7}
\]
The initial-segment relations, expanded with
\noderef{TreeSequenceInitialSegment}, and the vertex-set definition
\noderef{TreeSequenceVertexSet} put \(v_n\) in both \(V(T_n)\) and
\(V(T'_j)\).  Hence the hypotheses on \(P\), together with (7), are exactly
those of \noderef{ExtensionBaseCaseTwoObservation21}.  Choose its vertex
\(u\) and exact exit suffix \(Q=P[x,u]\).  In particular,
\[
 u\in V(T'_j)\setminus V(T_n),                                \tag{8}
\]
\(u\) is the sole vertex-node of \(Q\) in
\(V(T_n)\cup V(T'_j)\), and the edge by which \(Q\) leaves \(u\) has
\(\varrho_1\)-color \(\gamma_n\).

Fix \(\beta\in\overline{\varrho}_1(u)\).
By (8), \(u\notin V(T_n)\).  On the other hand,
\(v_n\) belongs to \(V(T_j)\), and the ordered inclusion
\(T_j\preccurlyeq T_n\), expanded with
\noderef{TreeSequenceInitialSegment} and
\noderef{TreeSequenceVertexSet}, puts it in \(V(T_n)\).  Thus \(u\ne v_n\).
Consequently
\(u\in V(T'_j)\setminus\{v_n\}\), and (C2) gives
\(\overline{\varrho}_1(u)\cap S'=\varnothing\).  The chosen membership
\(\beta\in\overline{\varrho}_1(u)\) therefore implies
\[
                              \beta\notin S'.                  \tag{9}
\]
Since \(\gamma_{1_1},\delta_n\in S'\), equation (9) gives
\(\beta\ne\gamma_{1_1}\) and \(\beta\ne\delta_n\).  Also \(u\ne v_n\),
and both vertices belong to the elementary set \(V(T'_j)\).  The
definition \noderef{ColoredVertexSetElementary}, applied to the
\(\beta\)-membership at \(u\), shows that \(\beta\) is not missing at
\(v_n\).  The ordered inclusions of the seed into \(T_j\) and \(T_n\),
together with (8), show that \(u\notin V(T_j(v_n))\): this is again the
vertex-set inclusion obtained from \noderef{TreeSequenceInitialSegment}
and \noderef{TreeSequenceVertexSet}.  Thus every
hypothesis of \noderef{ExtensionBaseCaseTwoObservation23} is available:
the b2 vertex and edge clauses and the displayed missing-set inclusion
are premises here.  Choose the resulting \(z,q_z\).  We retain
\[
 \beta\notin\overline{\varrho}_1(T_j(z)-z),\qquad
 \varrho_1(T_j.\operatorname{edge}(a))\ne\beta
 \quad(|T_j(v_n)|\le a<q_z).                                 \tag{10}
\]

We also specialize the already established generic
\noderef{ExtensionBaseCaseTwoObservation19}.  Its closure premise is the
given closure of \(T_j(v_n)\), its distinguished vertex is \(v_n\), and
its two missing-color premises are precisely (C2).  Therefore
\[
 \overline{\varrho}_1(T'_j)\cap S'=\{\delta_n\},               \tag{11}
\]
and no exact edge appended after \(T_j(v_n)\) in \(T'_j\) has a color in
\(S'\setminus\{\delta_n\}\).

Define \(\varrho_0\) as in (6).  The closure definition
\noderef{TashkinovClosure} says, in the sense of
\noderef{ColoredVertexSetClosed}, that \(V(T'_j)\) is
\(\varrho_1\)-closed.  Both swap colors are missing on the whole tree:
\(\beta\) is missing at \(u\in V(T'_j)\), and \(\delta_n\) is missing at
\(v_n\in V(T'_j)\), where the latter membership follows from the seed
initial segment by \noderef{TreeSequenceInitialSegment} and
\noderef{TreeSequenceVertexSet}.  The union definition
\noderef{TreeMissingColors} therefore puts both colors in
\(\overline{\varrho}_1(T'_j)\).  They are distinct by (9).  Hence
\noderef{OutsideTreeColorSwapStable} applies.  It proves that
\(\varrho_0\) is proper, leaves exactly \(e_0\) uncolored, is stable on
\(T'_j\), and, more strongly,
\[
 \varrho_0(a)=\varrho_1(a)
 \quad\text{for every exact identity incident with }V(T'_j). \tag{12}
\]

We next identify the component \(K\) switched in the definition of
\(\mu_0\).  By \noderef{TwoColorComponent}, the hypothesis that the support
of \(P\) is the full old component means that it contains every node
connected to \(v_n\) in the old two-color incidence graph.  At \(u\), the
path \(P\) has one edge toward \(x\), namely
the \(\gamma_n\)-colored exit identity from observation (21), and one
edge back toward \(v_n\), of color \(\delta_n\).  Indeed, the adjacency
clauses of \noderef{TwoColorIncidenceGraph} allow only these two colors,
and \noderef{IsProperEdgeColoring} excludes the already used
\(\gamma_n\) on the distinct adjacent edge.  The latter identity is
not \(f_n\), because observation (21) says that \(u\) is not incident
with \(f_n\).  Both identities are incident with \(T'_j\), so (12) fixes
them.  Past the exit identity, every vertex-node of the suffix \(Q-u\)
is outside \(V(T'_j)\); expansion of
\noderef{OutsideTreeColorSwap} therefore replaces every
\(\delta_n\)-edge of this alternating suffix by a \(\beta\)-edge and
leaves every \(\gamma_n\)-edge fixed.  Thus \(Q\) is a
\((\gamma_n,\beta)\)-path under \(\varrho_0\).  It cannot continue from
\(u\) toward \(v_n\), because that next old \(\delta_n\)-edge is fixed
and \(\beta\) is missing at \(u\).

The support of \(P\) is the full old component in the precise sense of
\noderef{TwoColorComponent}, and \(P\) is simple.  Every node of the old
incidence graph has degree at most two under the proper coloring, by
\noderef{TwoColorIncidenceDegreeAtMostTwo}.  The same holds after the
outside swap, which is proper.  If the
\((\gamma_n,\beta)\)-component through \(u\), as defined by
\noderef{TwoColorComponent}, had a node outside \(Q\), connectivity of
that component would give a shortest route from that node to \(Q\).  The
first new adjacency on such a route would either give an internal
node of \(Q\) a third neighbor or extend one of its endpoints.  The first
is impossible by the degree bound; the second is impossible at \(u\)
by the preceding paragraph and at \(x\) because the old full path ended
there and the outside transposition merely renamed
\(\delta_n\) to \(\beta\).  Consequently
\[
 K=\{y:y\in Q\},\qquad
 K\cap\{\mathsf{inl}(w):w\in V(T_n)\cup V(T'_j)\}
   =\{\mathsf{inl}(u)\}.                                     \tag{13}
\]
This proves the exact component-intersection assertion used below.

Since \(\gamma_n\ne\beta\) by (9) and
\(\varrho_0\) is proper,
\noderef{TwoColorComponentSwapProper} shows that \(\mu_0\) is proper.
Expansion of \noderef{TwoColorComponentSwap} shows pointwise that an
edge is uncolored under \(\mu_0\) exactly when it is uncolored under
\(\varrho_0\): values \(\gamma_n,\beta\) are exchanged on \(K\), and all
other option values are fixed.  Thus the two clauses of
\noderef{IsOneEdgeUncolored} show that \(\mu_0\) leaves exactly \(e_0\)
uncolored.

We prove (e1).  Every vertex of \(T_j(v_n)-v_n\) lies in
\(V(T_j)\subseteq V(T_n)\): the two inclusions follow from the ordered
prefix clauses of \noderef{TreeSequenceInitialSegment} and the vertex-set
description \noderef{TreeSequenceVertexSet}.  It is distinct from \(u\)
by (8).  Equation
(12) says that the first recoloring fixes every identity incident with
this set.  Equation (13), the adjacency definition in
\noderef{TwoColorIncidenceGraph}, the connected-component definition
\noderef{TwoColorComponent}, and the two endpoint identities in
\noderef{LooplessMultigraph} show that no edge-node of \(K\) is incident
with a vertex of this set; otherwise adjacency would put that vertex-node
in the same component \(K\).  The component switch therefore fixes all
those exact incident identities as well.  The incidence-witness definition of
\noderef{PresentColors} and complement definition of
\noderef{MissingColors} now give equality of the \(\mu_0\)- and
\(\varrho_1\)-missing sets at every preceding vertex.  Both clauses of
the b1 \noderef{StableColoringBeforeVertex} relation transfer unchanged
from \(\varrho_1\) to \(\mu_0\).  This is (e1).

At \(u\), equation (12) retains the fact that \(\beta\) is missing and
the exit identity supplies a present \(\gamma_n\).  The rooted switch
in \noderef{TwoColorComponentSwap} interchanges these two colors on every
incident edge-node in \(K\).  The incidence-witness and complement
definitions \noderef{PresentColors} and \noderef{MissingColors} therefore
show that \(\gamma_n\) becomes missing and \(\beta\) becomes present at
\(u\).  Combining this with the exact field identity (4) proves (e3).

For (e4), fix \(t\in L_1\setminus\{n\}\), with exact fields
\(\gamma_t,\delta_t\).  By (9) and the support-color premise, \(\beta\)
is different from both colors; the explicit color-identity premise also
makes \(\delta_n\) different from both.  Hence the outside transposition
does not create, remove, or recolor any \(\gamma_t\)- or
\(\delta_t\)-edge identity.  The corresponding two-color incidence
graph, its components, and the b4 strong-interchangeability certificates
are identical under \(\varrho_1\) and \(\varrho_0\).

The support-order premise puts every vertex-node of each cycle appearing
in those certificates in \(V(T_n)\).  By
\noderef{TwoColorComponent}, each such cycle support and \(K\) are supports
of connected components in their respective two-color incidence graphs,
while (13) says that \(K\) has no
vertex-node there.  If an edge-node belonged both to such an earlier
cycle and to \(K\), either endpoint, using \noderef{Incident} and the
adjacency clauses of \noderef{TwoColorIncidenceGraph}, would belong to both
components, contradicting their vertex-node disjointness.
Thus the component switch fixes every exact edge on every earlier cycle.
Conversely, no new earlier-color adjacency can enter those cycles:
an entering edge-node in \(K\) would again carry one of its incident
vertex-nodes into both components.  Their full component supports,
cycle witnesses, avoidance clauses, and support-order containments are
therefore unchanged.  Unfolding \noderef{ETTStrongInterchangeable}
proves (e4).

For (e5), equation (12) fixes every edge incident with \(T'_j\), and
(13) shows that the second switch fixes every edge incident with a vertex
of \(T'_j-u\).  Applying \noderef{PresentColors} and
\noderef{MissingColors} to these incident-edge equalities gives
\[
 \overline{\mu}_0(v)=\overline{\varrho}_1(v)
 \quad(v\in V(T'_j)\setminus\{u\}).                            \tag{14}
\]
The initial-segment clauses in \noderef{TreeSequenceInitialSegment}, read
through \noderef{TreeSequenceVertexSet}, put \(v_n\) in \(V(T'_j)\); since
\(u\ne v_n\), (14) holds there in particular.  Intersecting (14) with \(S'\) and
using (C2) and (11) gives the first two equalities of (e5).  Every
displayed edge identity of \(T'_j\) has both endpoints in \(V(T'_j)\).
It is fixed by (12), and if it belonged to \(K\), adjacency in
\noderef{TwoColorIncidenceGraph} and connectedness in
\noderef{TwoColorComponent} would put its endpoint different from \(u\)
in \(K\), contradicting (13); looplessness ensures that such an endpoint
exists.  It is therefore also fixed by the second switch.  The exact
appended-edge conclusion following (11) now gives the last clause of (e5).

It remains to verify (e2).  We first record the information on the
\(T_j\)-prefix.  Observation (23), equation (10), excludes \(\beta\) from
every vertex strictly before \(z\) and from every exact displayed edge
strictly between the seed and \(z\).  Since
\(\delta_n\notin\overline\varphi_{j-1}(T_j)\), the union definition
\noderef{TreeMissingColors} says that \(\delta_n\) is present under
\(\varphi_{j-1}\) at every vertex of \(T_j\).  At a vertex other than
\(v_n\), b2 transfers this presence to \(\varrho_1\).  We deliberately do
not make that transfer at \(v_n\), where (C2) instead says that
\(\delta_n\) is missing.

The global augmentation premise is a finite chain by
\noderef{TashkinovAugmentationSequence}.  At the unique one-step
augmentation which appends any exact position between the seed and \(z\),
\noderef{TashkinovAugmentation} says that its
\(\varphi_{j-1}\)-color is missing on the preceding tree.  It cannot be
\(\delta_n\), since \noderef{TreeMissingColors} excludes \(\delta_n\)
from every preceding-tree vertex.  The b2 edge equality transfers this
exact edge value to \(\varrho_1\); equation (10) separately excludes
\(\beta\).  Thus every such appended identity is fixed by the first swap.
Every seed identity is also an identity of \(T'_j\), by
\noderef{TreeSequenceInitialSegment}, and hence is fixed by (12).

Now split the missing-set assertion vertexwise.  If a stated vertex lies
in \(V(T'_j)\setminus\{u\}\), equation (14) gives its full missing-set
equality immediately.  This case includes \(v_n\); no presence of
\(\delta_n\) at \(v_n\) is asserted or needed.

Suppose instead that \(v\) is a stated vertex of \(T_j(z)\) outside
\(V(T'_j)\).  Since \(v\ne z\), observation (23) says that \(\beta\) is
present at \(v\).  Also \(v\ne v_n\), because the seed inclusion
\(T_j(v_n)\preccurlyeq T'_j\), expanded by
\noderef{TreeSequenceInitialSegment} and
\noderef{TreeSequenceVertexSet}, puts \(v_n\) in \(V(T'_j)\).  The first
paragraph of this verification therefore says that \(\delta_n\) is
present at \(v\) as well.

Neither of these two colors can occur on an edge from \(v\) into
\(V(T'_j)\).  Indeed, \noderef{TreeMissingColors} puts \(\beta\) and
\(\delta_n\) among the colors missing on \(T'_j\), witnessed respectively
by \(u\) and \(v_n\); \noderef{TashkinovClosure} supplies closedness, and
\noderef{ColoredVertexSetClosed} forbids either missing color on a boundary
identity.  By \noderef{BoundaryEdges}, an edge from the outside vertex
\(v\) into \(V(T'_j)\) is exactly such a boundary identity.  Consequently
every \(\beta\)- or \(\delta_n\)-edge incident with \(v\) is wholly outside
\(T'_j\), so \noderef{OutsideTreeColorSwap} transposes the two complete
nonempty incidence sets at \(v\).  Since both colors were present, this
transposition leaves the entire present set, and hence by
\noderef{PresentColors} and \noderef{MissingColors} the entire missing set,
at \(v\) unchanged.

The ordered inclusions \(T_j(z)\preccurlyeq T_j\preccurlyeq T_n\),
interpreted through \noderef{TreeSequenceInitialSegment} and
\noderef{TreeSequenceVertexSet}, give
\(V(T_j(z))\subseteq V(T_n)\).  Equation (13) and the component definition
\noderef{TwoColorComponent} therefore make the second switch disjoint from
every vertex of this prefix.  Thus the first-swap equality just proved is
also the required \(\mu_0\)-versus-\(\varrho_1\) equality.  Together with
the \(T'_j\)-case from (14), this proves the vertex equality in (e2).

The exact displayed edges are fixed as well.  We proved this for every
seed identity and every edge appended between the seed and \(z\); hence
the first swap fixes all of \(T_j(z)\).  Equation (12) fixes all of
\(T'_j(u)\).  For the second swap, the preceding containment and (13) show
that no edge of \(T_j(z)\) lies in \(K\): otherwise adjacency in
\noderef{TwoColorIncidenceGraph} and connectedness in
\noderef{TwoColorComponent} would put one of its endpoint vertices in
\(K\cap V(T_n)\).  Likewise, if a displayed edge of \(T'_j(u)\) lay in
\(K\), its two endpoint vertices, which lie in \(V(T'_j)\) by
\noderef{TreeSequenceVertexSet}, would lie in the connected component
\(K\) by those same two definitions; (13) permits only \(u\), and
looplessness rules out both endpoints being \(u\).  Thus the second swap
fixes every exact identity in both prefixes.

Finally suppose conditionally that \(\beta\) is missing at \(z\).  The
vertex \(z\) belongs to \(V(T_j)\subseteq V(T_n)\), by
\noderef{TreeSequenceInitialSegment} and
\noderef{TreeSequenceVertexSet}, whereas \(u\notin V(T_n)\); hence
\(z\ne u\).  Elementariness of \(V(T'_j)\), together with the simultaneous
\(\beta\)-membership at \(u\), now shows that \(z\notin V(T'_j)\).
Therefore \(z\ne v_n\), and the first paragraph gives that \(\delta_n\)
is present there.  The same closed-boundary argument using
\noderef{ColoredVertexSetClosed} and \noderef{BoundaryEdges} shows that
all incident \(\beta\)- and \(\delta_n\)-identities are outside \(T'_j\).
Expansion of \noderef{OutsideTreeColorSwap}, followed by
\noderef{PresentColors} and \noderef{MissingColors}, consequently shows
that the first transposition makes \(\delta_n\) missing at \(z\).
Since \(z\in V(T_n)\), (13) and \noderef{TwoColorComponent} show that the
second switch fixes every identity incident with \(z\), so
\(\delta_n\in\overline\mu_0(z)\).  This is the conditional clause, and all
three assertions of (e2) are proved.
\end{proof}
